\documentclass[11pt]{article}

\usepackage[margin=1in]{geometry}
\usepackage{amsmath,amssymb,amsthm}
\usepackage{graphicx}
\usepackage[hidelinks]{hyperref}
\usepackage{enumitem}
\usepackage{xcolor}

\setlist[itemize]{leftmargin=1.5em}
\setlist[enumerate]{leftmargin=1.7em}
\setlength{\parindent}{0pt}
\setlength{\parskip}{0.7em}

\newtheorem*{theorem}{Theorem}
\newtheorem*{claim}{Claim}
\newtheorem*{remark}{Remark}

\DeclareMathOperator{\supp}{supp}
\newcommand{\K}{\mathbb K}
\newcommand{\N}{\mathbb N}

\begin{document}

\begin{center}
{\Large Literature-Already-Answered Packet}\\[0.25em]
{\large The pair $(c_0,c_0)$ does not have the CPPm}\\[0.35em]
\textbf{Status:} literature already answered, not a new automatic-discovery result
\end{center}

\section*{Source and later answer}

The source paper is:
\[
\text{arXiv:2601.17316, Manwook Han, \emph{Weak minimizing property on pairs of classical Banach spaces}.}
\]
The local copy is included as \texttt{source\_paper.pdf}.

The later supporting paper is:
\[
\text{arXiv:2605.01397, Anselmo Raposo Jr. and Geivison Ribeiro,}
\]
\[
\text{\emph{Weak Minimizing Property and the Compact Perturbation Property for the Minimum Modulus}.}
\]
The local copy is included as \texttt{supporting\_paper\_2605.01397.pdf}.

\section*{Open question location}

The following full-width crop is taken from page 13 of the source PDF. It shows
Han's statement that the case of $(c_0,c_0)$ remained undecided and the exact
Question 3.6.

\begin{center}
\includegraphics[width=\textwidth]{figures/open_problem_crop.png}
\end{center}

\section*{Question}

Han asks:
\begin{quote}
Does the pair $(c_0,c_0)$ have the CPPm?
\end{quote}

Here CPPm is the compact perturbation property for the minimum modulus: for
every operator $T:X\to Y$ that does not attain its minimum modulus,
\[
\sup_{K\in \mathcal K(X,Y)} m(T+K)=m(T),
\qquad
m(T)=\inf\{\|Tx\|:x\in S_X\}.
\]

\section*{Answer}

Raposo--Ribeiro explicitly cite Han's Question 3.6 and answer it negatively.
The following is the constructive proof, rewritten in a self-contained form.

\begin{theorem}
The pair $(c_0,c_0)$ does not have the CPPm.
\end{theorem}

\begin{proof}
Let
\[
f(x)=\sum_{j=2}^{\infty}2^{1-j}x_j,\qquad x=(x_j)_{j\ge 1}\in c_0.
\]
Then $f\in (c_0)^*=\ell_1$, $\|f\|=1$, $f(e_1)=0$, and $f$ has infinite
support. Define
\[
T:c_0\to c_0,\qquad Tx=x-f(x)e_1.
\]

We first compute the minimum modulus of $T$. Let $x\in S_{c_0}$ and put
\[
s=\sup_{j\ge 2}|x_j|.
\]
Since $f$ only sees the coordinates $j\ge 2$, we have $|f(x)|\le s$, and
\[
\|Tx\|=\max\{|x_1-f(x)|,s\}.
\]
If $s\ge 1/2$, then $\|Tx\|\ge 1/2$. If $s<1/2$, then the equality
$\|x\|_\infty=1$ forces $|x_1|=1$, and therefore
\[
|x_1-f(x)|\ge |x_1|-|f(x)|\ge 1-s>1/2.
\]
Thus $\|Tx\|\ge 1/2$ for every $x\in S_{c_0}$.

For the reverse inequality, take
\[
x^{(N)}=e_1+\frac12\sum_{j=2}^{N}e_j.
\]
Then $\|x^{(N)}\|_\infty=1$ and
\[
f(x^{(N)})=\frac12\sum_{j=2}^{N}2^{1-j}\longrightarrow \frac12.
\]
Consequently
\[
\|Tx^{(N)}\|
=\max\left\{\left|1-f(x^{(N)})\right|,\frac12\right\}
\longrightarrow \frac12,
\]
so $m(T)=1/2$.

Next we show that $T$ does not attain its minimum modulus. Suppose, toward a
contradiction, that $\|x\|_\infty=1$ and $\|Tx\|=1/2$. Then $s\le 1/2$. Since
$\|x\|_\infty=1$, this implies $|x_1|=1$. Multiplying $x$ by a scalar of
modulus one if necessary, we may assume $x_1=1$.

The inequality $\|Tx\|=1/2$ gives
\[
|1-f(x)|\le \frac12,
\]
hence $|f(x)|\ge 1/2$. On the other hand,
\[
|f(x)|
\le \sum_{j=2}^{\infty}2^{1-j}|x_j|
\le \frac12\sum_{j=2}^{\infty}2^{1-j}
=\frac12.
\]
Therefore equality holds throughout. Since all weights $2^{1-j}$ are positive,
equality in the second inequality forces $|x_j|=1/2$ for every $j\ge 2$. This
contradicts $x\in c_0$. Hence $T$ is not minimum-modulus attaining.

Finally define the rank-one operator
\[
K:c_0\to c_0,\qquad Kx=f(x)e_1.
\]
It is compact, and
\[
(T+K)x=x\qquad (x\in c_0).
\]
Thus $T+K=I_{c_0}$, so
\[
m(T+K)=1>\frac12=m(T).
\]
This violates the defining CPPm equality for the non-minimum-modulus-attaining
operator $T$. Therefore $(c_0,c_0)$ does not have the CPPm.
\end{proof}

\section*{Classification}

This packet belongs in
\texttt{solutions/literature\_already\_answered/}. The later paper explicitly
states that Han left Question 3.6 open and that it gives a negative answer.
Therefore this is a successful retrieval and verification hit, not an original
automatic-discovery proof.

\section*{Verification code}

The script \texttt{code/verify\_cppm\_example.py} checks the finite truncations
of the minimizing sequence $x^{(N)}$ and verifies the perturbation identity
$T+K=I$ on sample finite vectors. The computation is not a substitute for the
analytic non-attainment argument; it is only a reproducibility sanity check for
the construction.

\section*{Reproducibility}

The screenshot crop was generated from the included source PDF by
\[
\texttt{code/crop\_open\_question.py}.
\]
All LaTeX build products are placed in \texttt{tmp/}. The final rendered packet
is \texttt{solution\_packet.pdf}.

\end{document}
